% Load the kaobook class
\documentclass[
	fontsize=10pt, % Base font size
	twoside=false, % Use different layouts for even and odd pages (in particular, if twoside=true, the margin column will be always on the outside)
	%open=any, % If twoside=true, uncomment this to force new chapters to start on any page, not only on right (odd) pages
%%	secnumdepth=1, % How deep to number headings. Defaults to 1 (sections)
]{kaobook}

\newcommand{\discipline}{Informatique -- PTSI}
\newcommand{\auteur}{Xavier Pessoles}
\newcommand{\repRel}{../../../..}

\input{\repRel/Style/packages_v2}
\newcommand{\repStyle}{\repRel/Style}
\input{\repRel/Style/macros_SII_v2}
\input{\repRel/Style/environment_v2}
\usepackage{\repStyle/UPSTI_pedagogique}


% Choose the language
\usepackage[english,french]{babel}
\frenchsetup{StandardItemLabels=true}



% Load the bibliography package
\usepackage{kaobiblio}
%\addbibresource{Cy_01_ModelisationSystemes.bib} % Bibliography file

% Load mathematical packages for theorems and related environments
\usepackage{kaotheorems}

% Load the package for hyperreferences
\usepackage{kaorefs}

\graphicspath{{images/}{./}} % Paths where images are looked for

\makeindex[columns=3, title=Alphabetical Index, intoc] % Make LaTeX produce the files required to compile the index

\newcommand{\zonecarree}[2]{%
\begin{center}
\begin{tikzpicture}
  \draw[step=5mm, lightgray, very thin] (0,0) grid (#1,#2);
  \draw[thick] (0,0) rectangle (#1,#2);
\end{tikzpicture}
\end{center}
}

\begin{document}

%%----------------------------------------------------------------------------------------
%%	BOOK INFORMATION
%%----------------------------------------------------------------------------------------

\titlehead{Document Template}
\title[SII en PSI* - Xavier Pessoles]{SII en PSI* - Xavier Pessoles}
\author[XP]{Xavier Pessoles}
\date{\today}
%\publishers{An Awesome Publisher}

%----------------------------------------------------------------------------------------
%
\frontmatter % Denotes the start of the pre-document content, uses roman numerals


%%----------------------------------------------------------------------------------------
%%	MAIN BODY
%%----------------------------------------------------------------------------------------
%
\mainmatter % Denotes the start of the main document content, resets page numbering and uses arabic 

\setchapterstyle{kao}

\setcounter{margintocdepth}{\sectiontocdepth}
%\marginlayout % pour etre en 2/3 1/3
%\fullwidthlayout

\graphicspath{{\repStyle/png}}

\pagestyle{xp.scrheadings}



%%%%%
\newcommand{\repExo}{}
\newcommand{\nomExo}{}
\AtBeginEnvironment{corrige}{\small}

\livrettrue % Livrettrue :  eleve sans les corrections
\collefalse\colletrue
\normaltrue
%%%%%

\newcommand{\exer}{\subsection*}



%%%%%%%%%%%%%%%%%%%%%%%%%%%%%%%%%%%%%%%%%
\renewcommand{\repExo}{\repRel/PSI_Cy_03_ConceptionCommande/Chapitre_01_Correction//}
\renewcommand{\nomExo}{Cy_03_01_TD_PI_03_SysReeduc}

\graphicspath{{\repRel/Style/png/}{\repExo/\nomExo/images/}}

\chapter*{Les interros qui ne comptent pas  -- 02}

\question{Ecrire la fonction \lstinline{moyenne(L:[int]) -> int} qui renvoie la moyenne des éléments de la liste \lstinline{L}.}
\ifprof
\else
\zonecarree{17}{3.5} 
\fi


\question{Ecrire la fonction \lstinline{maxi(L:[int]) -> int} qui renvoie la valeur du maximum de la liste \lstinline{L}.}

\ifprof
\else
\zonecarree{17}{3} 
\fi


\question{Ecrire la fonction \lstinline{indice_max(L:[int]) -> int} qui renvoie l'indice de la valeur du maximum de la liste \lstinline{L}.}

\ifprof
\else
\zonecarree{17}{3} 
\fi

\question{Ecrire la fonction \lstinline{mini(L:[int]) -> int} qui renvoie la valeur du minimum de la liste \lstinline{L}.}

\ifprof
\else
\zonecarree{17}{3} 
\fi


\question{Ecrire la fonction \lstinline{indice_min(L:[int]) -> int} qui renvoie l'indice de la valeur du minimum de la liste \lstinline{L}.}

\ifprof
\else
\zonecarree{17}{3} 
\fi

\question{Ecrire la fonction \lstinline{pairs(L:[int]) -> [int]} qui renvoie la liste des éléments pairs de la liste \lstinline{L}.}

\ifprof
\else
\zonecarree{17}{3} 
\fi

\question{Ecrire la fonction \lstinline{del_max(L:[int]) -> [int]} qui permet de renvoyer une liste identique à \lstinline{L} sans son maximum.}

\ifprof
\else
\zonecarree{17}{3} 
\fi

\question{Ecrire la fonction \lstinline{compte(L:[int], x:int) -> int} qui permet de compter le nombre d'occurences de \lstinline{x} dans la liste \lstinline{L}.}

\ifprof
\else
\zonecarree{17}{3} 
\fi



\end{document}

